À \qty{20}{\degreeCelsius}, la solubilité d'un sel dans l'eau est de \qty{\v{S}}{g} par \qty{100}{mL} d'eau. On ajoute \qty{\v{m}}{g} de ce sel à \qty{\v{Ve}}{mL} d'eau à \qty{20}{\degreeCelsius} et on agite longtemps.

\begin{enumerate}[label=\alph*)]
    \item Quelle masse maximale de sel peut se dissoudre dans cette eau ?
    \item Quelle masse de sel reste au fond du contenant ? Comment appelle-t-on la solution obtenue ?
    \item Quel volume minimal d'eau à \qty{20}{\degreeCelsius} faudrait-il pour dissoudre tout le sel ?
\end{enumerate}
